\subsection{Objectives and Scope}
\label{sec:objectives}

The research presented in this thesis is inspired by and limited in scope to map reuse in episodic KV-SLAM operations. Limiting the scope means that some relevant topics will not be explored in depth, but will be mentioned in Section \ref{sec:future_work} as potential areas for follow-on research. The justification for the scope selection is given here.

The methods presented in this research will not work for all SLAM systems, or on all SLAM-generated maps. While change detection techniques do exist for dense maps \cite{PLACEHOLDERa}, the method presented here would not be suitable, as it relies on locating the same map features at different points in time. Dense maps identify matching positions through geometric correspondence, as opposed to feature descriptors. This applies to LiDAR-based SLAM systems as well, as no descriptor is generated for each point. The method presented relies on the capability to apply and persist metadata to a particular map point, something which is not possible without point descriptors, leaving KV-SLAM as the only option for this work.

The focus on map reuse and the episodic operating model can be motivated by the use cases discussed previously. On the ISS, the ideal navigation system would be a self-updating localizer rather than a SLAM system. The primary difference is that a localizer utilizes the same map and reference frame for all operations, while a SLAM system generates and uses its own map. SLAM's features are extremely useful for exploration of unknown areas, but less so for repeated operations in known areas. The combination of map reuse and the episodic operating model allows the system to act as a self-updating localizer, maintaining a map that is up to date with the last observed state of the world, while preserving as much previously collected data as possible. Within this scope, we can define the objectives of this research.

This research has two primary objectives. The first is the implementation of a novel Map Point Removal (MPR) method that incorporates historical, viewpoint-dependent observability at the point level using a K-Nearest Neighbors (KNN) Observability Model. The second is the integration of this MPR method into an existing KV-SLAM system. Both objectives have several sub-goals, which are detailed below.

\subsubsection*{Objective 1: Implementation of a Viewpoint-Dependent Map Point Removal (MPR) Method}
\label{sec:objective_1}
\subparagraph{Goal 1: Speed}

As this system is intended to be run as part of real-time SLAM operations, no function performed by the model can take longer than the delta time between image frames. Ideally, each function should be significantly faster than this to allow for all SLAM system operations to complete as well.

This goal also refers to the speed of identifying outdated map points, meaning a point that is outdated should be identified and culled with a minimal number of observations.

\subparagraph{Goal 2: Size}

To run well on resource-constrained hardware, additional space requirements (RAM and disk) must be kept low. The space complexity of the system should be constant, and not a function of the number of map points or observations.

\subparagraph{Goal 3: Accuracy}

The system must only identify points as outdated once sufficient evidence of them being out of date has been collected. The system should be robust to falsely identifying active map points as outdated and outdated map points as active.

\subsubsection*{Objective 2: Integration into an Existing KV-SLAM System}

To verify the performance and effectiveness of the proposed MPR method, it must be implemented into an existing SLAM system and benchmarked in various scenarios.

\subparagraph{Goal 1: Maintain Performance in the General Case}

This research is focused on outdated map reuse, and therefore should not affect the performance of the SLAM system when not reusing an outdated map. The modified system should perform identically to the base system in this case.

\subparagraph{Goal 2: Improve Performance When Using Outdated Maps}

To be considered successful, the SLAM implementation should perform better than the base system when using an outdated initial map. The criteria for better performance are defined in Section \ref{sec:eval_metrics}.

\subparagraph{Goal 3: Maintain Computational Requirements}

If the goals of Objective 1 are met, integration of the MPR method should not cause the computational requirements of the system to increase significantly.

\bigskip
\noindent With these objectives in hand, we can now define a set of questions which, when answered, will determine whether the goals of the objectives have been successfully achieved.
